\chapter{Lezione 16 --- Osservatore asintotico, dualità e principio di separazione}

\section{Richiamo: osservatore dello stato}

Consideriamo il sistema lineare tempo-invariante
\[
\begin{cases}
\dot x = Ax + Bu,\\[4pt]
y = Cx,
\end{cases}
\]
dove:
\[
x(t)\in\mathbb{R}^n,\qquad
u(t)\in\mathbb{R}^m,\qquad
y(t)\in\mathbb{R}^\ell,
\]
e
\[
A\in\mathbb{R}^{n\times n},\qquad
B\in\mathbb{R}^{n\times m},\qquad
C\in\mathbb{R}^{\ell\times n}.
\]

Lo stato \(x(t)\) non è direttamente accessibile. Per questo si introduce uno
\textbf{stato osservato} o \textbf{stimato}
\[
\hat x(t),
\]
costruito mediante un osservatore dinamico.

\section{Equazioni dell'osservatore}

L'osservatore considerato a lezione è del tipo di Luenberger:
\[
\begin{cases}
\dot{\hat x} = A\hat x + Bu + L(y-\hat y),\\[4pt]
\hat y = C\hat x,
\end{cases}
\]
dove
\[
L\in\mathbb{R}^{n\times \ell}
\]
è la matrice di guadagno dell'osservatore.

L'idea è semplice:
\begin{itemize}
    \item copio la dinamica del sistema;
    \item confronto l'uscita vera \(y\) con l'uscita stimata \(\hat y\);
    \item correggo la stima mediante il termine \(L(y-\hat y)\).
\end{itemize}

\section{Errore di osservazione}

Definiamo l'errore di osservazione:
\[
e := x-\hat x.
\]

Allora
\[
\dot e = \dot x - \dot{\hat x}.
\]

Sostituendo le equazioni del sistema e dell'osservatore:
\[
\dot e
=
(Ax+Bu)-\bigl(A\hat x+Bu+L(y-\hat y)\bigr).
\]

Poiché
\[
y-\hat y = Cx-C\hat x = C(x-\hat x)=Ce,
\]
si ottiene
\[
\dot e = A(x-\hat x)-LC(x-\hat x),
\]
cioè
\[
\dot e = (A-LC)e.
\]

Questa è la \textbf{dinamica dell'errore di osservazione}.

\section{Osservatore asintotico}

Se tutti gli autovalori della matrice
\[
A-LC
\]
hanno parte reale strettamente negativa, allora
\[
e(t)\to 0 \qquad \text{per } t\to+\infty.
\]

Quindi
\[
\hat x(t)\to x(t).
\]

Per questo motivo l'equazione dell'osservatore viene chiamata
\textbf{osservatore asintotico}.

\subsection*{Interpretazione}

Non imponiamo che \(\hat x(t)\) coincida subito con \(x(t)\),
ma facciamo in modo che converga asintoticamente allo stato vero.

\section{Scelta di \(L\) e assegnamento degli autovalori dell'osservatore}

La matrice \(L\) è scelta da noi.  
Quindi il problema è del tutto analogo a quello già visto per la retroazione di stato:
\begin{itemize}
    \item lì si progettava \(K\) per modificare la matrice \(A-BK\);
    \item qui si progetta \(L\) per modificare la matrice \(A-LC\).
\end{itemize}

Vale allora il seguente risultato.

\begin{theorem}
Se la coppia \((A,C)\) è completamente osservabile, allora esiste una matrice
\[
L\in\mathbb{R}^{n\times \ell}
\]
tale da assegnare arbitrariamente gli autovalori della matrice
\[
A-LC.
\]
\end{theorem}

Quindi, sotto ipotesi di osservabilità completa, possiamo scegliere la velocità
di convergenza dell'osservatore.

\section{Dualità tra osservabilità e controllabilità}

Negli appunti si sottolinea che \textbf{osservabilità e controllabilità sono concetti duali}.

Dato il sistema
\[
\begin{cases}
\dot x = Ax + Bu,\\
y = Cx,
\end{cases}
\]
il suo sistema duale è
\[
\begin{cases}
\dot z = A^T z + C^T w,\\[4pt]
\rho = B^T z.
\end{cases}
\]

La proprietà chiave è:
\[
(A,C)\ \text{completamente osservabile}
\quad\Longleftrightarrow\quad
(A^T,C^T)\ \text{completamente controllabile}.
\]

Allora, applicando al sistema duale il teorema di assegnamento degli autovalori
mediante retroazione completa dello stato, sappiamo che esiste una matrice \(K_d\)
tale che gli autovalori di
\[
A^T-C^T K_d
\]
siano arbitrari.

Ma
\[
A^T-C^T K_d = (A-K_d^T C)^T,
\]
e una matrice ha gli stessi autovalori della sua trasposta. Dunque
\[
\lambda(A^T-C^T K_d)=\lambda(A-K_d^T C).
\]

Questo mostra che possiamo porre
\[
L = K_d^T,
\]
e assegnare arbitrariamente gli autovalori di
\[
A-LC.
\]

Questa è la giustificazione teorica del progetto dell'osservatore tramite dualità.

\section{Schema a blocchi dell'osservatore}

Dal punto di vista strutturale:
\begin{itemize}
    \item il sistema vero evolve con stato \(x\) e uscita \(y\);
    \item l'osservatore evolve con stato \(\hat x\) e uscita \(\hat y=C\hat x\);
    \item la differenza \(y-\hat y\) viene amplificata da \(L\) e reiniettata nell'osservatore.
\end{itemize}

Un punto concettuale importante è che lo \textbf{stato osservato} \(\hat x\)
è sempre accessibile, perché l'osservatore è un algoritmo che costruiamo noi.
Quindi, anche se lo stato vero \(x\) non è misurabile, lo stato \(\hat x\) lo è.

\section{Retroazione sullo stato osservato}

Poiché \(\hat x\) è disponibile, possiamo usarlo nella legge di controllo:
\[
u=-K\hat x.
\]

Questa è la soluzione pratica al problema della non accessibilità dello stato:
invece di retroazionare lo stato vero, retroazioniamo il suo stimatore.

\section{Equazioni del sistema accoppiato}

Consideriamo quindi il sistema vero
\[
\begin{cases}
\dot x = Ax + Bu,\\
y = Cx,
\end{cases}
\]
e l'osservatore
\[
\begin{cases}
\dot{\hat x} = A\hat x + Bu + L(y-\hat y),\\
\hat y = C\hat x,
\end{cases}
\]
con legge di controllo
\[
u=-K\hat x.
\]

\subsection{Dinamica del sistema vero}

Sostituendo \(u=-K\hat x\) nella dinamica del sistema:
\[
\dot x = Ax - BK\hat x.
\]

Poiché
\[
\hat x = x-e,
\]
segue
\[
\dot x = A x - BK(x-e)
= (A-BK)x + BK\,e.
\]

\subsection{Dinamica dell'osservatore}

L'equazione dell'osservatore diventa
\[
\dot{\hat x}
=
A\hat x - BK\hat x + L(Cx-C\hat x),
\]
cioè
\[
\dot{\hat x}
=
(A-BK-LC)\hat x + LC\,x.
\]

\subsection{Dinamica dell'errore}

Calcoliamo ora di nuovo l'errore:
\[
e=x-\hat x.
\]

Derivando:
\[
\dot e = \dot x - \dot{\hat x}.
\]

Sostituendo le espressioni precedenti:
\[
\dot e
=
\bigl((A-BK)x + BK e\bigr)
-
\bigl((A-BK)\hat x + LC(x-\hat x)\bigr).
\]

Poiché \(x-\hat x=e\), si ottiene
\[
\dot e = (A-BK)e + BK e - LC e,
\]
e quindi
\[
\dot e = (A-LC)e.
\]

Questo è un risultato fondamentale:
la dinamica dell'errore di osservazione \textbf{non dipende da \(K\)},
ma solo da \(L\).

\section{Dinamica congiunta di stato ed errore}

Se come nuovo vettore di stato prendiamo
\[
\begin{bmatrix}
x\\
e
\end{bmatrix},
\]
la dinamica complessiva diventa
\[
\begin{bmatrix}
\dot x\\
\dot e
\end{bmatrix}
=
\underbrace{
\begin{bmatrix}
A-BK & BK\\
0 & A-LC
\end{bmatrix}
}_{D}
\begin{bmatrix}
x\\
e
\end{bmatrix}.
\]

La matrice \(D\) è \textbf{triangolare a blocchi}.
Quindi il suo polinomio caratteristico è
\[
p_D(\lambda)
=
\det(\lambda I-D)
=
\det\!\bigl(\lambda I-(A-BK)\bigr)\,
\det\!\bigl(\lambda I-(A-LC)\bigr).
\]

Ne segue che gli autovalori di \(D\) sono l'unione di:
\begin{itemize}
    \item autovalori di \(A-BK\), legati alla dinamica del controllo;
    \item autovalori di \(A-LC\), legati alla dinamica dell'osservatore.
\end{itemize}

\section{Principio di separazione}

Il risultato precedente è il \textbf{principio di separazione}.

\begin{quote}
Gli autovalori del controllore (\(A-BK\)) e quelli dell'osservatore (\(A-LC\))
sono assegnabili \emph{indipendentemente}.
\end{quote}

Questo significa che:
\begin{itemize}
    \item posso progettare \(K\) come se avessi accesso diretto allo stato;
    \item posso progettare \(L\) come se stessi costruendo soltanto l'osservatore;
    \item il sistema complessivo ha come autovalori l'unione dei due insiemi.
\end{itemize}

\subsection*{Attenzione}

Il principio di separazione \textbf{non} significa che la dinamica effettiva di \(x\)
sia identica al caso in cui si retroaziona direttamente lo stato vero.

Infatti nella dinamica di \(x\) compare il termine
\[
BK\,e,
\]
che agisce come perturbazione finché l'errore di osservazione non è ancora nullo.

Quindi:
\begin{itemize}
    \item i \emph{modi} del sistema sono assegnabili separatamente;
    \item ma lo stato vero \(x\) è influenzato dall'errore \(e\) finché questo non si è estinto.
\end{itemize}

Se però
\[
e(t)\to 0,
\]
allora asintoticamente la retroazione sullo stato osservato si comporta come
la retroazione sullo stato vero.

\section{Scelta degli autovalori dell'osservatore}

Poiché vogliamo che l'errore di osservazione si annulli rapidamente,
conviene scegliere gli autovalori dell'osservatore più veloci di quelli del sistema controllato.

Nel caso continuo questo significa tipicamente scegliere gli autovalori di
\[
A-LC
\]
più a sinistra nel piano complesso rispetto a quelli di
\[
A-BK.
\]

In altre parole:
\begin{itemize}
    \item il sistema vero deve essere stabilizzato e avere la dinamica desiderata;
    \item l'osservatore deve convergere ancora più rapidamente,
    così l'errore \(e\) si annulla prima che domini la dinamica del sistema.
\end{itemize}

\subsection*{Osservazione pratica}

Negli appunti compare una nota importante:
\begin{itemize}
    \item sul sistema fisico non posso rendere arbitrariamente veloci gli autovalori, perché rischierei
    di richiedere segnali di controllo troppo grandi e saturare gli attuatori;
    \item per l'osservatore questo problema è meno critico, perché è un algoritmo interno.
\end{itemize}

Quindi in genere si scelgono gli autovalori dell'osservatore significativamente più veloci
di quelli del controllore.

\section{Interpretazione strutturale}

Dal punto di vista concettuale, il sistema completo è costituito da due sottosistemi:
\begin{enumerate}
    \item il \textbf{sistema vero}, con stato \(x\);
    \item il \textbf{sistema osservato}, con stato \(\hat x\).
\end{enumerate}

L'osservatore utilizza:
\begin{itemize}
    \item l'ingresso \(u\), che è noto;
    \item l'uscita misurata \(y\), che è nota;
    \item il proprio stato interno \(\hat x\), che è accessibile.
\end{itemize}

Una volta che \(\hat x\) converge a \(x\), posso trattare la retroazione
\[
u=-K\hat x
\]
come un sostituto pratico di
\[
u=-Kx.
\]

\section{Sintesi della lezione}

I punti principali della lezione sono:
\begin{enumerate}
    \item l'osservatore asintotico è descritto da
    \[
    \dot{\hat x}=A\hat x+Bu+L(y-\hat y),
    \qquad
    \hat y=C\hat x;
    \]

    \item definendo
    \[
    e=x-\hat x,
    \]
    la dinamica dell'errore è
    \[
    \dot e=(A-LC)e;
    \]

    \item se tutti gli autovalori di \(A-LC\) hanno parte reale negativa,
    allora
    \[
    e(t)\to 0
    \quad\text{e quindi}\quad
    \hat x(t)\to x(t);
    \]

    \item se la coppia \((A,C)\) è completamente osservabile, gli autovalori di \(A-LC\)
    si possono assegnare arbitrariamente;

    \item questo risultato deriva per dualità dalla teoria della retroazione di stato;

    \item usando la retroazione sullo stato osservato
    \[
    u=-K\hat x,
    \]
    il sistema vero soddisfa
    \[
    \dot x=(A-BK)x + BK\,e;
    \]

    \item la dinamica complessiva di \((x,e)\) è
    \[
    \begin{bmatrix}
    \dot x\\
    \dot e
    \end{bmatrix}
    =
    \begin{bmatrix}
    A-BK & BK\\
    0 & A-LC
    \end{bmatrix}
    \begin{bmatrix}
    x\\
    e
    \end{bmatrix};
    \]

    \item gli autovalori del sistema complessivo sono l'unione di quelli di \(A-BK\)
    e di quelli di \(A-LC\);

    \item controllore e osservatore si progettano indipendentemente:
    questo è il principio di separazione;

    \item in pratica conviene scegliere gli autovalori dell'osservatore più veloci
    di quelli del sistema controllato.
\end{enumerate}